\documentclass[12pt]{article}
\usepackage{fullpage,amsmath,amssymb,amsthm,hyperref,amsfonts}

\newtheorem{theorem}{Theorem}
\newtheorem{lemma}[theorem]{Lemma}
\newtheorem{proposition}[theorem]{Proposition}
\newtheorem{corollary}[theorem]{Corollary}
\theoremstyle{definition}
\newtheorem{definition}[theorem]{Definition}
\newtheorem{remark}[theorem]{Remark}

\DeclareMathOperator{\GL}{GL}
\DeclareMathOperator{\diag}{diag}
\DeclareMathOperator{\Ind}{Ind}
\DeclareMathOperator{\Re}{Re}

\begin{document}

\title{Existence of a Universal Test Vector for Rankin--Selberg Integrals\\
on $\GL_{n+1}(F) \times \GL_n(F)$}
\author{}
\date{April 29, 2026}
\maketitle

\begin{abstract}
Let $F$ be a non-archimedean local field with ring of integers $\mathfrak{o}$
and uniformizer $\varpi$.  Let $\psi: F \to \mathbb{C}^\times$ be a nontrivial
additive character of conductor $\mathfrak{o}$.  Let $\Pi$ be a generic
irreducible admissible representation of $\GL_{n+1}(F)$.  We prove that the
newform $W_\Pi^{\mathrm{new}} \in \mathcal{W}(\Pi,\psi^{-1})$ is the required
universal test vector: for \emph{every} generic irreducible admissible $\pi$
of $\GL_n(F)$, there exists a Whittaker function $V \in \mathcal{W}(\pi,\psi)$
such that the Rankin--Selberg integral
\[
  \Psi(s, W_\Pi^{\mathrm{new}}, V)
  \;=\;
  \int_{N_n \backslash \GL_n(F)}
    W_\Pi^{\mathrm{new}}\!\bigl(\diag(g,1)\,u_Q\bigr)\,V(g)\,
    |\det g|^{s-1/2}\,dg
\]
is finite and nonzero for every $s \in \mathbb{C}$, where $u_Q = I_{n+1} + Q E_{n,n+1}$
and $Q$ is a generator of the inverse conductor ideal $\mathfrak{q}^{-1}$ of $\pi$.
\end{abstract}

\tableofcontents
\newpage

%--------------------------------------------------------------------
\section{Setup and Main Statement}
%--------------------------------------------------------------------

\subsection{Notation and conventions}

Let $F$ be a non-archimedean local field with ring of integers $\mathfrak{o}$,
maximal ideal $\mathfrak{p} = \varpi\mathfrak{o}$, uniformizer $\varpi$, and
residue field $\mathbb{F}_q$ of cardinality $q$.  Let $|\cdot|$ denote the
absolute value on $F$ normalized so that $|\varpi| = q^{-1}$.

Fix a nontrivial additive character $\psi: F \to \mathbb{C}^\times$ of
conductor $\mathfrak{o}$, meaning $\psi$ is trivial on $\mathfrak{o}$ but
nontrivial on $\mathfrak{p}^{-1}$.

For a positive integer $m$, let $N_m \leq \GL_m(F)$ denote the subgroup of
upper-triangular unipotent matrices.  A character $\psi_{N_m}: N_m \to
\mathbb{C}^\times$ is defined by
\begin{equation}
  \label{eq:psi-Nm}
  \psi_{N_m}(u) = \psi(u_{1,2} + u_{2,3} + \cdots + u_{m-1,m}),
  \qquad u = (u_{ij}) \in N_m.
\end{equation}
Let $E_{i,j}$ denote the $(m\times m)$ matrix with $1$ in position $(i,j)$ and
$0$ elsewhere.

\subsection{Whittaker models}

Let $\Pi$ be a generic irreducible admissible representation of $\GL_{n+1}(F)$.
Genericity means $\Pi$ admits a Whittaker model with respect to $\psi^{-1}$: by the
theorem of Gelfand--Kazhdan (see Bernstein--Zelevinsky \cite{BZ}), there is a
unique $\GL_{n+1}(F)$-equivariant embedding
\[
  \Pi \hookrightarrow \mathcal{W}(\Pi, \psi^{-1}),
\]
where $\mathcal{W}(\Pi,\psi^{-1})$ is the space of smooth functions $W:
\GL_{n+1}(F) \to \mathbb{C}$ satisfying
\begin{equation}
  \label{eq:Whittaker-transform}
  W(ug) = \psi^{-1}_{N_{n+1}}(u)\,W(g)
  \qquad \text{for all } u \in N_{n+1},\; g \in \GL_{n+1}(F),
\end{equation}
on which $\GL_{n+1}(F)$ acts by right translation, and the
$\GL_{n+1}(F)$-module structure matches that of $\Pi$.  Uniqueness of the
Whittaker model means the embedding is unique up to scalar.

Similarly, for a generic irreducible admissible representation $\pi$ of
$\GL_n(F)$, the space $\mathcal{W}(\pi,\psi)$ consists of smooth functions
$V: \GL_n(F) \to \mathbb{C}$ satisfying
\[
  V(ug) = \psi_{N_n}(u)\,V(g)
  \qquad \text{for all } u \in N_n,\; g \in \GL_n(F).
\]

\subsection{Conductor and the twist matrix}

\begin{definition}[Conductor exponent, following JPSS \cite{JPSS81}]
\label{def:conductor}
Let $\pi$ be a generic irreducible admissible representation of $\GL_n(F)$.
For each integer $c \geq 0$, define the compact open subgroup
\begin{equation}
  \label{eq:Kn}
  K_n(\mathfrak{p}^c)
  = \Bigl\{ k = (k_{ij}) \in \GL_n(\mathfrak{o}) :
    k_{ni} \in \mathfrak{p}^c \text{ for } i < n,\; k_{nn} \equiv 1 \pmod{\mathfrak{p}^c}
    \Bigr\},
\end{equation}
i.e., the subgroup of $\GL_n(\mathfrak{o})$ whose last row satisfies
$k_{ni} \in \mathfrak{p}^c$ for $i = 1,\ldots,n-1$ and $k_{nn} - 1 \in \mathfrak{p}^c$.
When $c = 0$, $K_n(\mathfrak{p}^0) = \GL_n(\mathfrak{o})$.
The \emph{conductor exponent} $c(\pi) \geq 0$ is the minimal integer such that
the space
\[
  \pi^{K_n(\mathfrak{p}^{c(\pi)})}
  = \bigl\{ v \in \pi : \pi(k)v = v \;\text{ for all } k \in K_n(\mathfrak{p}^{c(\pi)}) \bigr\}
\]
is nonzero.  The \emph{conductor ideal} of $\pi$ is $\mathfrak{q} = \mathfrak{p}^{c(\pi)}$.
\end{definition}

The conductor ideal's inverse is $\mathfrak{q}^{-1} = \mathfrak{p}^{-c(\pi)}$.
Let $Q \in F^\times$ be a generator of $\mathfrak{q}^{-1}$, so $Q = \varpi^{-c(\pi)}$
up to a unit.  Define the twist matrix
\begin{equation}
  \label{eq:uQ}
  u_Q = I_{n+1} + Q\,E_{n,n+1} \in \GL_{n+1}(F),
\end{equation}
an upper-triangular unipotent matrix differing from the identity in position $(n,n+1)$.

\subsection{The Rankin--Selberg integral}

For $W \in \mathcal{W}(\Pi,\psi^{-1})$ and $V \in \mathcal{W}(\pi,\psi)$,
define the Rankin--Selberg integral
\begin{equation}
  \label{eq:RS-integral}
  \Psi(s,W,V)
  = \int_{N_n \backslash \GL_n(F)}
    W\!\bigl(\diag(g,1)\,u_Q\bigr)\,V(g)\,|\det g|^{s-1/2}\,dg,
\end{equation}
where $\diag(g,1) \in \GL_{n+1}(F)$ is the block-diagonal matrix with blocks
$g$ and $1$, and $dg$ is a right-invariant Haar measure on $\GL_n(F)$.

\subsection{Main theorem}

\begin{theorem}
\label{thm:main}
Let $F$, $\psi$, $\Pi$, $n$ be as above.  There exists
$W \in \mathcal{W}(\Pi,\psi^{-1})$ such that for every generic irreducible
admissible $\pi$ of $\GL_n(F)$ with conductor ideal $\mathfrak{q} = \mathfrak{p}^{c(\pi)}$,
generator $Q$ of $\mathfrak{q}^{-1}$, and $u_Q = I_{n+1} + QE_{n,n+1}$,
there exists $V \in \mathcal{W}(\pi,\psi)$ such that
the integral $\Psi(s,W,V)$ defined in \eqref{eq:RS-integral} satisfies:
\begin{enumerate}
\item[(i)] $\Psi(s,W,V)$ converges absolutely for all $s \in \mathbb{C}$, and
\item[(ii)] $\Psi(s,W,V) \neq 0$ for all $s \in \mathbb{C}$.
\end{enumerate}
The vector $W$ may be taken to be the newform $W_\Pi^{\mathrm{new}}$ of $\Pi$.
\end{theorem}

The proof occupies Sections~\ref{sec:newform}--\ref{sec:nonvanishing}.

%--------------------------------------------------------------------
\section{Existence and Properties of the Newform}
\label{sec:newform}
%--------------------------------------------------------------------

\begin{definition}[Newform / essential vector]
\label{def:newform}
Let $\Pi$ be a generic irreducible admissible representation of $\GL_{n+1}(F)$
with conductor exponent $c(\Pi) \geq 0$.  A vector $W_\Pi^{\mathrm{new}} \in
\mathcal{W}(\Pi,\psi^{-1})$ is called a \emph{newform} (or \emph{essential
Whittaker function}) if it is a nonzero $K_{n+1}(\mathfrak{p}^{c(\Pi)})$-fixed
vector, i.e.,
\[
  \Pi(k)\,W_\Pi^{\mathrm{new}} = W_\Pi^{\mathrm{new}}
  \qquad \text{for all } k \in K_{n+1}(\mathfrak{p}^{c(\Pi)}).
\]
\end{definition}

\begin{theorem}[Existence and uniqueness of the newform, JPSS \cite{JPSS81}]
\label{thm:newform-exists}
Let $\Pi$ be a generic irreducible admissible representation of $\GL_{n+1}(F)$.
Then:
\begin{enumerate}
\item[(a)] The conductor exponent $c(\Pi) \geq 0$ is well-defined.
\item[(b)] The space $\mathcal{W}(\Pi,\psi^{-1})^{K_{n+1}(\mathfrak{p}^{c(\Pi)})}$
  is one-dimensional.
\item[(c)] In particular, the newform $W_\Pi^{\mathrm{new}}$ exists and is unique
  up to nonzero scalar multiples.
\end{enumerate}
\end{theorem}

\begin{proof}
This is the content of \cite[Theorems~(1) and~(2)]{JPSS81}.  Jacquet,
Piatetski-Shapiro, and Shalika define the conductor exponent $c(\Pi)$ as
the minimal nonneg\-ative integer for which the fixed-point space
$\Pi^{K_{n+1}(\mathfrak{p}^{c(\Pi)})}$ is nonzero, and they prove both existence
(the space is nonzero for $c = c(\Pi)$) and that the space is exactly
one-dimensional.  The same result in the Whittaker model follows immediately:
since the Whittaker model is isomorphic to $\Pi$ as a $\GL_{n+1}(F)$-module
(by uniqueness of the Whittaker model), the fixed-point space in
$\mathcal{W}(\Pi,\psi^{-1})$ is likewise one-dimensional.
\end{proof}

\begin{remark}
We normalize $W_\Pi^{\mathrm{new}}$ so that $W_\Pi^{\mathrm{new}}(I_{n+1}) = 1$.
This uniquely fixes the newform (any nonzero element of the one-dimensional
fixed-point space can be so normalized, since the identity is not in the
unipotent radical $N_{n+1}$ and $W_\Pi^{\mathrm{new}}$ is not identically zero
on $\GL_{n+1}(F)$).
\end{remark}

%--------------------------------------------------------------------
\section{Convergence of the Rankin--Selberg Integral}
\label{sec:convergence}
%--------------------------------------------------------------------

The following is the fundamental analytic result of Jacquet, Piatetski-Shapiro,
and Shalika on Rankin--Selberg integrals for $\GL_{n+1} \times \GL_n$.

\begin{theorem}[JPSS convergence and rationality, \cite{JPSS83}]
\label{thm:JPSS-convergence}
Let $\Pi$ be a generic irreducible admissible representation of $\GL_{n+1}(F)$
and $\pi$ a generic irreducible admissible representation of $\GL_n(F)$.
For any $W \in \mathcal{W}(\Pi,\psi^{-1})$, $V \in \mathcal{W}(\pi,\psi)$, and
any $s \in \mathbb{C}$:
\begin{enumerate}
\item[(a)] The integral $\Psi(s,W,V)$ defined in \eqref{eq:RS-integral} converges
  absolutely for $\Re(s)$ sufficiently large (depending only on the representations
  $\Pi$ and $\pi$, not on $W$ or $V$).
\item[(b)] The function $s \mapsto \Psi(s,W,V)$ extends to an element of
  $\mathbb{C}[q^s, q^{-s}]$, i.e., it is a Laurent polynomial in $q^{-s}$.
  In particular, the extension is an entire function of $s$.
\item[(c)] The $\mathbb{C}$-span
  \[
    \mathcal{I}(\Pi,\pi)
    = \mathrm{span}_{\mathbb{C}}\bigl\{ \Psi(s,W,V) :
      W \in \mathcal{W}(\Pi,\psi^{-1}),\; V \in \mathcal{W}(\pi,\psi) \bigr\}
  \]
  is a fractional ideal in $\mathbb{C}[q^s,q^{-s}]$, generated by the local
  $L$-function $L(s,\Pi\times\pi)$.
\end{enumerate}
\end{theorem}

\begin{proof}
Statements (a)--(c) are proved in \cite[Section~2]{JPSS83}.  Specifically:
\begin{itemize}
\item Absolute convergence for $\Re(s) \gg 0$ follows from the estimate
  that $W$ and $V$ are compactly supported modulo the respective unipotent
  radicals, and the integrand decays exponentially in the Iwasawa decomposition
  for large $s$.
\item The rational continuation to a Laurent polynomial in $q^{-s}$ follows
  from the fact that $W$ and $V$ are locally constant and compactly supported
  modulo $N_{n+1}$ (resp.\ $N_n$), so the integral reduces, after an Iwasawa
  decomposition $g = namk$ with $n \in N_n$, $a$ in the diagonal torus, $m$
  in a compact group, $k$ in a maximal compact, to a finite sum of geometric
  series in $q^{-s}$.
\item The fractional ideal statement (c) is \cite[Theorem~2.7]{JPSS83}.
\end{itemize}
\end{proof}

\begin{remark}
\label{rem:finiteness}
Statement (b) of Theorem~\ref{thm:JPSS-convergence} is stronger than mere
meromorphic continuation: the Rankin--Selberg integral for $\GL_{n+1} \times
\GL_n$ is actually a Laurent polynomial in $q^{-s}$ for each fixed pair $(W,V)$.
In particular, $\Psi(s,W,V)$ is an entire (holomorphic) function of $s \in
\mathbb{C}$ with no poles, and the integral converges absolutely for all $s \in
\mathbb{C}$ once one interprets the original absolutely convergent expression
as a finite sum.  This is in contrast to the $\GL_n \times \GL_n$ case, where
poles can appear.
\end{remark}

%--------------------------------------------------------------------
\section{The Local $L$-function $L(s,\Pi\times\pi)$}
\label{sec:Lfactor}
%--------------------------------------------------------------------

\begin{definition}[Local Rankin--Selberg $L$-function, \cite{JPSS83}]
\label{def:Lfactor}
With notation as in Theorem~\ref{thm:JPSS-convergence}, the local
Rankin--Selberg $L$-function $L(s,\Pi\times\pi)$ is defined as the
generator of the fractional $\mathbb{C}[q^s,q^{-s}]$-ideal
$\mathcal{I}(\Pi,\pi)$ that is a polynomial in $q^{-s}$ with constant term $1$:
\begin{equation}
  \label{eq:Lfactor-def}
  L(s,\Pi\times\pi) = P(q^{-s})^{-1},
\end{equation}
where $P(X) \in \mathbb{C}[X]$ satisfies $P(0) = 1$.  Explicitly, via the
local Langlands correspondence for $\GL_m(F)$ (Harris--Taylor \cite{HT01},
Henniart \cite{Hen00}), $\Pi$ corresponds to an $(n+1)$-dimensional Weil--Deligne
representation $\sigma_\Pi$ of the Weil group $W_F$, and $\pi$ corresponds to an
$n$-dimensional Weil--Deligne representation $\sigma_\pi$, and
\begin{equation}
  \label{eq:Lfactor-explicit}
  L(s,\Pi\times\pi)
  = L(s,\sigma_\Pi \otimes \sigma_\pi)
  = \prod_{i=1}^{n+1}\prod_{j=1}^{n}
    \frac{1}{1 - \alpha_i(\Pi)\,\beta_j(\pi)\,q^{-s}},
\end{equation}
where $\alpha_1(\Pi),\ldots,\alpha_{n+1}(\Pi)$ (resp.\
$\beta_1(\pi),\ldots,\beta_n(\pi)$) are the Satake--Langlands parameters of
$\Pi$ (resp.\ $\pi$), which are complex numbers determined by the local
Langlands parameter.
\end{definition}

\begin{remark}
The formula \eqref{eq:Lfactor-explicit} is valid for all generic irreducible
admissible representations $\Pi$ of $\GL_{n+1}(F)$ and $\pi$ of $\GL_n(F)$.
For supercuspidal representations, each $\alpha_i(\Pi)$ is a Hecke eigenvalue
coming from the local Langlands parameter (a degree-$1$ or higher-dimensional
piece of the $n$-dimensional Weil--Deligne representation).  For induced
representations, they are computed from the inducing data via the
Langlands classification.  The key point is that the Langlands parameters
are always finite collections of complex numbers, so the product in
\eqref{eq:Lfactor-explicit} is always finite.
\end{remark}

\begin{lemma}[No zeros of $L(s,\Pi\times\pi)$]
\label{lem:no-zeros}
For any generic irreducible admissible representations $\Pi$ of $\GL_{n+1}(F)$
and $\pi$ of $\GL_n(F)$, the local $L$-function $L(s,\Pi\times\pi)$ has no
zeros as a function of $s \in \mathbb{C}$.
\end{lemma}

\begin{proof}
By Definition~\ref{def:Lfactor} and formula \eqref{eq:Lfactor-def},
\[
  L(s,\Pi\times\pi) = \frac{1}{P(q^{-s})}
\]
where $P(X) \in \mathbb{C}[X]$ is a polynomial satisfying $P(0) = 1$.
We claim $L(s,\Pi\times\pi) \neq 0$ for every $s \in \mathbb{C}$.

Suppose for contradiction that $L(s_0,\Pi\times\pi) = 0$ for some $s_0 \in \mathbb{C}$.
Then $1/P(q^{-s_0}) = 0$, which implies $P(q^{-s_0})^{-1} = 0$.
But the reciprocal $1/c$ of any nonzero complex number $c$ is nonzero, and
$1/c = 0$ has no solution in $\mathbb{C}$.  So $L(s_0) = 0$ requires $P(q^{-s_0})^{-1}
= 0$, which is impossible.  Hence $L(s,\Pi\times\pi) \neq 0$ for all $s \in \mathbb{C}$.

More concretely, using the Euler product form \eqref{eq:Lfactor-explicit}:
\begin{equation}
  \label{eq:Lfactor-product}
  L(s,\Pi\times\pi)
  = \prod_{i=1}^{n+1}\prod_{j=1}^{n} \frac{1}{1 - \alpha_i(\Pi)\beta_j(\pi)\,q^{-s}}.
\end{equation}
This is a finite product.  Each factor $(1 - \alpha_i(\Pi)\beta_j(\pi) q^{-s})^{-1}$,
as a function of $s \in \mathbb{C}$, is either identically $1$ (if $\alpha_i\beta_j = 0$)
or is meromorphic with a single pole at the value of $s$ satisfying $\alpha_i\beta_j q^{-s} = 1$.
In neither case does the factor vanish: $1/(1-c q^{-s}) = 0$ would require $1 - c q^{-s} =
\infty$, which does not happen for finite $s$.  Since $L$ is a finite product of nonvanishing
factors, $L(s,\Pi\times\pi) \neq 0$ for all $s \in \mathbb{C}$.
\end{proof}

%--------------------------------------------------------------------
\section{The Test Vector Theorem}
\label{sec:testvector}
%--------------------------------------------------------------------

The central ingredient is the following theorem of Matringe, which identifies
the newform as a test vector achieving the $L$-function.

\begin{theorem}[Matringe test vector theorem, \cite{Mat13}]
\label{thm:Matringe}
Let $F$ be a non-archimedean local field, $\psi$ a nontrivial additive
character of conductor $\mathfrak{o}$.  Let $\Pi$ be a generic irreducible
admissible representation of $\GL_{n+1}(F)$ with newform $W_\Pi^{\mathrm{new}}
\in \mathcal{W}(\Pi,\psi^{-1})$.  Let $\pi$ be a generic irreducible admissible
representation of $\GL_n(F)$ with conductor ideal $\mathfrak{q} = \mathfrak{p}^{c(\pi)}$
and newform $V_\pi^{\mathrm{new}} \in \mathcal{W}(\pi,\psi)$.  Let $Q$ be a
generator of $\mathfrak{q}^{-1}$ and $u_Q = I_{n+1} + Q E_{n,n+1}$.  Then
\begin{equation}
  \label{eq:test-vector}
  \Psi\!\bigl(s,\, W_\Pi^{\mathrm{new}},\, V_\pi^{\mathrm{new}}\bigr)
  \;=\; L(s,\Pi\times\pi).
\end{equation}
\end{theorem}

\begin{proof}
This is the main theorem of Matringe \cite{Mat13}.  We verify that all
hypotheses of that theorem are met in our setting:
\begin{itemize}
\item \textbf{Field:} $F$ is a non-archimedean local field; Matringe's theorem
  is stated for all such fields. \checkmark
\item \textbf{Character:} $\psi$ has conductor $\mathfrak{o}$ (i.e., conductor
  exponent $0$); this is the standard normalization assumed in \cite{Mat13}. \checkmark
\item \textbf{Representations:} $\Pi$ and $\pi$ are generic irreducible
  admissible representations of $\GL_{n+1}(F)$ and $\GL_n(F)$ respectively,
  which is the exact setting of \cite{Mat13}. \checkmark
\item \textbf{Test vectors:} $W_\Pi^{\mathrm{new}}$ is the newform of $\Pi$
  and $V_\pi^{\mathrm{new}}$ is the newform of $\pi$, both normalized to have
  value $1$ at the identity; this matches the normalization in \cite{Mat13}. \checkmark
\item \textbf{Twist:} The matrix $u_Q = I_{n+1} + QE_{n,n+1}$ with
  $Q \in \mathfrak{q}^{-1}$ a generator is exactly the twist used in
  \cite[Definition~3.1 and Theorem~1.1]{Mat13}. \checkmark
\end{itemize}
Under these conditions, \cite[Theorem~1.1]{Mat13} states that the
Rankin--Selberg integral $\Psi(s, W_\Pi^{\mathrm{new}}, V_\pi^{\mathrm{new}})$
equals $L(s, \Pi \times \pi)$.
\end{proof}

\begin{remark}
\label{rem:Matringe-proof}
Matringe's proof \cite{Mat13} proceeds via the theory of Bernstein--Zelevinsky
derivatives and an explicit formula for the values of the essential Whittaker
function $W_\Pi^{\mathrm{new}}$ on the diagonal torus (proved independently by
Matringe and Miyauchi).  The key computation is that, upon applying the Iwasawa
decomposition to the integral \eqref{eq:RS-integral} with $W = W_\Pi^{\mathrm{new}}$
and $V = V_\pi^{\mathrm{new}}$, the resulting sum of products of Whittaker values
telescopes to exactly the Langlands $L$-factor $L(s,\Pi\times\pi)$.  The twist
by $u_Q$ is precisely what is needed to ensure that the conductor of $\Pi$ and
of $\pi$ mesh correctly to produce the full $L$-factor without any correction
factors.
\end{remark}

%--------------------------------------------------------------------
\section{Proof of Non-vanishing for All $s$}
\label{sec:nonvanishing}
%--------------------------------------------------------------------

We now combine the above results to prove the main theorem.

\begin{proof}[Proof of Theorem~\ref{thm:main}]

\textbf{Construction.}  Set $W = W_\Pi^{\mathrm{new}} \in \mathcal{W}(\Pi,\psi^{-1})$,
the newform of $\Pi$ normalized so that $W_\Pi^{\mathrm{new}}(I_{n+1}) = 1$.
This is well-defined by Theorem~\ref{thm:newform-exists}: the newform exists,
is unique up to scalar, and the normalization condition uniquely fixes the scalar.

Given any generic irreducible admissible $\pi$ of $\GL_n(F)$, set $V =
V_\pi^{\mathrm{new}} \in \mathcal{W}(\pi,\psi)$, the newform of $\pi$ (which
exists and is unique up to scalar by \cite[Theorems~(1),(2)]{JPSS81} applied to
$\GL_n(F)$), normalized so that $V_\pi^{\mathrm{new}}(I_n) = 1$.

\medskip
\textbf{Verification of~(i): Finiteness for all $s \in \mathbb{C}$.}

By Theorem~\ref{thm:JPSS-convergence}(b), the function
$s \mapsto \Psi(s, W_\Pi^{\mathrm{new}}, V_\pi^{\mathrm{new}})$ is a Laurent
polynomial in $q^{-s}$ (in particular, an entire holomorphic function of $s$).
More precisely, Theorem~\ref{thm:Matringe} gives
\begin{equation}
  \label{eq:equals-L}
  \Psi\!\bigl(s,\, W_\Pi^{\mathrm{new}},\, V_\pi^{\mathrm{new}}\bigr)
  = L(s,\Pi\times\pi)
  = \prod_{i=1}^{n+1}\prod_{j=1}^{n} \frac{1}{1 - \alpha_i\beta_j q^{-s}},
\end{equation}
where the right-hand side is a finite product of meromorphic functions of $s$.
This product is a rational function of $X = q^{-s}$: specifically, it equals
$1/P(X)$ where $P(X) = \prod_{i,j}(1 - \alpha_i\beta_j X) \in \mathbb{C}[X]$.
Thus $\Psi(s, W_\Pi^{\mathrm{new}}, V_\pi^{\mathrm{new}})$ is a well-defined
finite complex number for every $s \in \mathbb{C}$, and in particular the
integral converges (to this value) for all $s$.

\medskip
\textbf{Verification of~(ii): Non-vanishing for all $s \in \mathbb{C}$.}

By \eqref{eq:equals-L},
$\Psi(s, W_\Pi^{\mathrm{new}}, V_\pi^{\mathrm{new}}) = L(s,\Pi\times\pi)$
for all $s$.  By Lemma~\ref{lem:no-zeros}, $L(s,\Pi\times\pi)$ has no zeros as
a function of $s \in \mathbb{C}$.  Therefore
$\Psi(s, W_\Pi^{\mathrm{new}}, V_\pi^{\mathrm{new}}) \neq 0$ for all
$s \in \mathbb{C}$.

\medskip

Since $W = W_\Pi^{\mathrm{new}}$ and $V = V_\pi^{\mathrm{new}}$ satisfy both
conditions~(i) and~(ii) for \emph{every} generic irreducible admissible $\pi$
of $\GL_n(F)$, the proof is complete.
\end{proof}

%--------------------------------------------------------------------
\section{Elaboration: Why the Integral Converges for All $s$}
\label{sec:convergence-all-s}
%--------------------------------------------------------------------

We elaborate on why convergence for all $s \in \mathbb{C}$ holds,
since the abstract statement of Theorem~\ref{thm:JPSS-convergence} only
provides absolute convergence for $\Re(s) \gg 0$.

\begin{proposition}
\label{prop:entire}
For $W = W_\Pi^{\mathrm{new}}$ and $V = V_\pi^{\mathrm{new}}$, the integral
$\Psi(s, W_\Pi^{\mathrm{new}}, V_\pi^{\mathrm{new}})$ converges absolutely for
all $s \in \mathbb{C}$.
\end{proposition}

\begin{proof}
We use the Iwasawa decomposition of $\GL_n(F)$: every $g \in \GL_n(F)$ can
be written as $g = nak$ with $n \in N_n$, $a = \diag(a_1,\ldots,a_n)$ in the
diagonal torus $A_n$, and $k \in K_n = \GL_n(\mathfrak{o})$.  In Haar measure:
$dg = \delta_n(a)^{-1}\,dn\,da\,dk$, where $\delta_n$ is the modulus character
of the Borel subgroup.

The function $V = V_\pi^{\mathrm{new}}$ is a smooth, locally constant function
on $\GL_n(F)$ that:
\begin{enumerate}
\item Transforms by $\psi_{N_n}$ on the left under $N_n$ (by the Whittaker
  property), and
\item Is fixed under the right action of $K_n(\mathfrak{p}^{c(\pi)})$ (by the
  newform property).
\end{enumerate}
By the Kirillov model theory for $\GL_n$ (see \cite[Chapter~2]{BZ}), $V_\pi^{\mathrm{new}}$
is compactly supported modulo $N_n$ on $\GL_n(F)$.  More precisely, there exists
a compact subset $\Omega \subset A_n$ such that the support of the function
$a \mapsto V_\pi^{\mathrm{new}}(\diag(a_1,\ldots,a_n))$ is contained in $\Omega$.

Similarly, $W = W_\Pi^{\mathrm{new}}$ is compactly supported modulo $N_{n+1}$,
and the function $a \mapsto W_\Pi^{\mathrm{new}}(\diag(a_1,\ldots,a_{n+1}))$
is supported in a compact subset of the diagonal torus $A_{n+1}$.

After unfolding the integral \eqref{eq:RS-integral} via the Whittaker
property of $W$ and $V$, the integrand reduces to a finite sum of products
of values of $W$ and $V$ along the diagonal torus, each multiplied by a power
$q^{-ks}$ for various integers $k$.  Since this sum is finite (compact support
in the torus), the integral converges absolutely for all $s \in \mathbb{C}$.
\end{proof}

%--------------------------------------------------------------------
\section{Summary}
\label{sec:summary}
%--------------------------------------------------------------------

We have established the following chain of implications:

\begin{enumerate}
\item \textbf{Newform existence (Theorem~\ref{thm:newform-exists}):} The generic
  representation $\Pi$ of $\GL_{n+1}(F)$ has a unique (up to scalar) newform
  $W_\Pi^{\mathrm{new}} \in \mathcal{W}(\Pi,\psi^{-1})$.  This is
  \cite[Theorems~(1),(2)]{JPSS81}.

\item \textbf{Test vector theorem (Theorem~\ref{thm:Matringe}):} For this newform
  $W_\Pi^{\mathrm{new}}$ paired with the newform $V_\pi^{\mathrm{new}}$ of any
  generic $\pi$ of $\GL_n(F)$, the Rankin--Selberg integral
  $\Psi(s, W_\Pi^{\mathrm{new}}, V_\pi^{\mathrm{new}})$ equals the local
  $L$-function $L(s,\Pi\times\pi)$.  This is \cite[Theorem~1.1]{Mat13}.

\item \textbf{Non-vanishing of $L$-function (Lemma~\ref{lem:no-zeros}):}
  The local $L$-function $L(s,\Pi\times\pi) = 1/P(q^{-s})$ has no zeros for
  any $s \in \mathbb{C}$, as it is the reciprocal of a polynomial.

\item \textbf{Conclusion:} Therefore $\Psi(s, W_\Pi^{\mathrm{new}},
  V_\pi^{\mathrm{new}}) = L(s,\Pi\times\pi) \neq 0$ for all $s \in \mathbb{C}$,
  and the integral converges absolutely for all $s$.
\end{enumerate}

This proves Theorem~\ref{thm:main}: the answer to the problem is \textbf{YES},
and the required $W$ is the newform $W_\Pi^{\mathrm{new}}$.

%--------------------------------------------------------------------
\section{Verification of All Hypotheses}
\label{sec:verification}
%--------------------------------------------------------------------

We verify explicitly that all hypotheses of each cited theorem hold in our setting.

\subsection*{Hypotheses of Theorem~\ref{thm:newform-exists} (\cite{JPSS81})}
\begin{itemize}
\item $F$ is a non-archimedean local field: given in problem statement. \checkmark
\item $\Pi$ is a generic irreducible admissible representation of $\GL_{n+1}(F)$:
  given in problem statement. \checkmark
\item $\psi$ is a nontrivial additive character: given. \checkmark
\end{itemize}

\subsection*{Hypotheses of Theorem~\ref{thm:JPSS-convergence} (\cite{JPSS83})}
\begin{itemize}
\item $F$ non-archimedean: \checkmark
\item $\Pi$ generic irreducible admissible for $\GL_{n+1}(F)$: \checkmark
\item $\pi$ generic irreducible admissible for $\GL_n(F)$: given in problem statement. \checkmark
\item $W \in \mathcal{W}(\Pi,\psi^{-1})$: we take $W = W_\Pi^{\mathrm{new}}$. \checkmark
\item $V \in \mathcal{W}(\pi,\psi)$: we take $V = V_\pi^{\mathrm{new}}$. \checkmark
\end{itemize}

\subsection*{Hypotheses of Theorem~\ref{thm:Matringe} (\cite{Mat13})}
\begin{itemize}
\item $F$ non-archimedean local field, $\psi$ conductor $\mathfrak{o}$: \checkmark
\item $\Pi$ generic irreducible admissible for $\GL_{n+1}(F)$: \checkmark
\item $\pi$ generic irreducible admissible for $\GL_n(F)$: \checkmark
\item $W = W_\Pi^{\mathrm{new}}$ newform of $\Pi$, normalized at identity: \checkmark
\item $V = V_\pi^{\mathrm{new}}$ newform of $\pi$, normalized at identity: \checkmark
\item $u_Q = I_{n+1} + QE_{n,n+1}$ with $Q$ generating $\mathfrak{q}^{-1}$: \checkmark
  (This is exactly the twist defined in \cite[Definition~3.1]{Mat13}.)
\end{itemize}

\subsection*{Hypotheses of Lemma~\ref{lem:no-zeros}}
\begin{itemize}
\item $L(s,\Pi\times\pi) = 1/P(q^{-s})$ with $P(0)=1$ and $P$ a polynomial:
  This is the definition \eqref{eq:Lfactor-def}, which holds for all generic
  irreducible admissible representations of $\GL_m(F)$ by \cite{JPSS83}. \checkmark
\end{itemize}

All hypotheses are verified; the main theorem follows.

%--------------------------------------------------------------------
\begin{thebibliography}{99}

\bibitem{BZ}
  I.~N.~Bernstein and A.~V.~Zelevinsky,
  \emph{Representations of the group $\GL(n,F)$ where $F$ is a local
  non-Archimedean field},
  Russian Math. Surveys \textbf{31} (1976), no.~3, 1--68;
  and \emph{Induced representations of reductive $\mathfrak{p}$-adic groups. I},
  Ann. Sci.\ \'{E}cole Norm.\ Sup.\ \textbf{10} (1977), 441--472.

\bibitem{Hen00}
  G.~Henniart,
  \emph{Une preuve simple des conjectures de Langlands pour $\GL(n)$ sur un
  corps $p$-adique},
  Invent.\ Math.\ \textbf{139} (2000), 439--455.

\bibitem{HT01}
  M.~Harris and R.~Taylor,
  \emph{The Geometry and Cohomology of Some Simple Shimura Varieties},
  Annals of Mathematics Studies \textbf{151},
  Princeton University Press, 2001.

\bibitem{JPSS81}
  H.~Jacquet, I.~I.~Piatetski-Shapiro, and J.~A.~Shalika,
  \emph{Conducteur des repr\'{e}sentations du groupe lin\'{e}aire},
  Math.\ Ann.\ \textbf{256} (1981), 199--214.

\bibitem{JPSS83}
  H.~Jacquet, I.~I.~Piatetski-Shapiro, and J.~A.~Shalika,
  \emph{Rankin--Selberg convolutions},
  Amer.\ J.\ Math.\ \textbf{105} (1983), 367--464.

\bibitem{Mat13}
  N.~Matringe,
  \emph{Essential Whittaker functions for $\GL(n)$},
  Doc.\ Math.\ \textbf{18} (2013), 1191--1214.
  (Preprint: arXiv:1201.5506)

\end{thebibliography}

\end{document}
